\chapter{\IfLanguageName{dutch}{Stand van zaken}{State of the art}}%
\label{ch:stand-van-zaken}

% Tip: Begin elk hoofdstuk met een paragraaf inleiding die beschrijft hoe
% dit hoofdstuk past binnen het geheel van de bachelorproef. Geef in het
% bijzonder aan wat de link is met het vorige en volgende hoofdstuk.

% Pas na deze inleidende paragraaf komt de eerste sectiehoofding.

%Dit hoofdstuk bevat je literatuurstudie. De inhoud gaat verder op de inleiding, maar zal het onderwerp van de bachelorproef *diepgaand* uitspitten. De bedoeling is dat de lezer na lezing van dit hoofdstuk helemaal op de hoogte is van de huidige stand van zaken (state-of-the-art) in het onderzoeksdomein. Iemand die niet vertrouwd is met het onderwerp, weet nu voldoende om de rest van het verhaal te kunnen volgen, zonder dat die er nog andere informatie moet over opzoeken \autocite{Pollefliet2011}.
%
%Je verwijst bij elke bewering die je doet, vakterm die je introduceert, enz.\ naar je bronnen. In \LaTeX{} kan dat met het commando \texttt{$\backslash${textcite\{\}}} of \texttt{$\backslash${autocite\{\}}}. Als argument van het commando geef je de ``sleutel'' van een ``record'' in een bibliografische databank in het Bib\LaTeX{}-formaat (een tekstbestand). Als je expliciet naar de auteur verwijst in de zin (narratieve referentie), gebruik je \texttt{$\backslash${}textcite\{\}}. Soms is de auteursnaam niet expliciet een onderdeel van de zin, dan gebruik je \texttt{$\backslash${}autocite\{\}} (referentie tussen haakjes). Dit gebruik je bv.~bij een citaat, of om in het bijschrift van een overgenomen afbeelding, broncode, tabel, enz. te verwijzen naar de bron. In de volgende paragraaf een voorbeeld van elk.
%
%\textcite{Knuth1998} schreef een van de standaardwerken over sorteer- en zoekalgoritmen. Experten zijn het erover eens dat cloud computing een interessante opportuniteit vormen, zowel voor gebruikers als voor dienstverleners op vlak van informatietechnologie~\autocite{Creeger2009}.
%
%Let er ook op: het \texttt{cite}-commando voor de punt, dus binnen de zin. Je verwijst meteen naar een bron in de eerste zin die erop gebaseerd is, dus niet pas op het einde van een paragraaf.
%
%\begin{figure}
%  \centering
%  \includegraphics[width=0.8\textwidth]{grail.jpg}
%  \caption[Voorbeeld figuur.]{\label{fig:grail}Voorbeeld van invoegen van een figuur. Zorg altijd voor een uitgebreid bijschrift dat de figuur volledig beschrijft zonder in de tekst te moeten gaan zoeken. Vergeet ook je bronvermelding niet!}
%\end{figure}
%
%\begin{listing}
%  \begin{minted}{python}
%    import pandas as pd
%    import seaborn as sns
%
%    penguins = sns.load_dataset('penguins')
%    sns.relplot(data=penguins, x="flipper_length_mm", y="bill_length_mm", hue="species")
%  \end{minted}
%  \caption[Voorbeeld codefragment]{Voorbeeld van het invoegen van een codefragment.}
%\end{listing}
%
%\lipsum[7-20]
%
%\begin{table}
%  \centering
%  \begin{tabular}{lcr}
%    \toprule
%    \textbf{Kolom 1} & \textbf{Kolom 2} & \textbf{Kolom 3} \\
%    $\alpha$         & $\beta$          & $\gamma$         \\
%    \midrule
%    A                & 10.230           & a                \\
%    B                & 45.678           & b                \\
%    C                & 99.987           & c                \\
%    \bottomrule
%  \end{tabular}
%  \caption[Voorbeeld tabel]{\label{tab:example}Voorbeeld van een tabel.}
%\end{table}

\subsection{RustProv}

RustProv is de naam van de provisioner die in deze bachelorproef wordt ontwikkeld. De naam combineert \emph{Rust}, de gekozen programmeertaal, en \emph{Prov}, een afkorting van provisioning.




\subsection{Virtualisatie en provisioning}

Virtuele machines vormen de basis van de omgevingen waarvoor de provisioner wordt ontwikkeld. In deze bachelorproef ligt de focus op het automatisch aanmaken, configureren en beschikbaar maken van een VM in een onderwijscontext. Virtualisatie laat toe om meerdere geïsoleerde omgevingen op dezelfde fysieke machine te gebruiken, wat handig is voor labo’s, testen en demo-omgevingen.

\subsection{Waarom een provisioner?}

Een provisioner automatiseert het opzetten en configureren van virtuele omgevingen. Daardoor verlopen testen en deployments consistenter, sneller en met minder kans op fouten. In deze bachelorproef is een provisioner vooral nuttig om een stabiele en herhaalbare omgeving te voorzien voor onderwijs- en labo-opdrachten. Dit vermijdt ook dat dezelfde configuratie telkens manueel opnieuw moet worden uitgevoerd.

\subsection{Infrastructure as Code}

Het automatiseringsproces in systeembeheer wordt vaak aangeduid als \emph{Infrastructure as Code (IaC)}. Binnen IaC wordt een onderscheid gemaakt tussen verschillende soorten tools, zoals \emph{provisioning tools}, \emph{configuration management} en \emph{orchestration}. Provisioningtools richten zich op het creëren van de infrastructuur, terwijl configuratiemanagementtools instaan voor de installatie en configuratie van software op de server. Het voordeel hiervan is dat de opbouw van een omgeving reproduceerbaar en beter beheersbaar wordt.

Voor lokale ontwikkelomgevingen worden vaak tools zoals Vagrant gebruikt. Daarnaast bestaan er tools zoals Ansible en Terraform, die eerder geassocieerd worden met cloudomgevingen of grotere configuraties. Welke tool gekozen wordt, hangt af van de context, de grootte van de omgeving en de mate van automatisering die nodig is.

\subsection{Educatieve context}

Hoewel er verschillende commerciële oplossingen bestaan voor enterprise-omgevingen, gelden in een onderwijscontext andere eisen voor virtualisatie. Onderzoek naar het gebruik van virtuele machines in het onderwijs toont dat een evenwicht nodig is tussen performantie, gebruiksvriendelijkheid voor studenten en de beperkingen van consumentenhardware zoals laptops. In een labo-omgeving moet een virtuele machine dus niet alleen werken, maar ook snel beschikbaar zijn.

Ook speelt tijd een grote rol. De \emph{provisioning time}, de tijd tussen het aanvragen van de virtuele machine en de beschikbaarheid ervan, is een kritieke performance-indicator die de productiviteit beïnvloedt. Lange provisioning times kunnen tijdens labo’s of examens leiden tot minder effectieve werktijd. Daardoor blijft de vraag bestaan hoe een lichtgewicht en snelle provisioner gerealiseerd kan worden die specifiek afgestemd is op Windows- en Linuxomgevingen in een educatieve context. Dit is extra relevant wanneer meerdere studenten tegelijk met dezelfde tool moeten werken.

\subsection{Rust als programmeertaal}

Rust werd gekozen als programmeertaal voor de provisioner vanwege de goede cross-platform compatibiliteit, waardoor de tool zowel op Linux als op Windows kan draaien zonder grote aanpassingen. De taal combineert performantie met geheugenveiligheid via het ownership-model, waardoor veelvoorkomende fouten zoals geheugenlekken worden vermeden. Dit is belangrijk voor infrastructuurtools, omdat fouten zoals een onderbroken provisioning of een mislukte SSH-verbinding gecontroleerd moeten kunnen worden afgehandeld. Rust is relatief jong, maar wordt steeds vaker gebruikt voor systeemsoftware en tools waar betrouwbaarheid belangrijk is.

\subsection{Afbakening van de doelomgeving}

De provisioner wordt ontworpen voor de concrete SEP-opdracht en moet de stappen ondersteunen die nodig zijn om die opdracht in een vooraf opgezette virtuele omgeving uit te voeren. Dit maakt de tool doelgericht en beperkt de scope tot wat effectief nodig is voor de opleiding. Door die afbakening blijft het project haalbaar binnen de beschikbare tijd en blijft de focus liggen op een werkende proof-of-concept. Dit project gebruikt een bridge adapter voor het netwerk en drie test-VM’s die eerder voor deze opdracht werden ingezet.


\chapter{Provisioning}
\label{ch:Provisioning}

\subsection{Tijdswinsten tijdens het provisioningsproces}

\subsubsection{VDI-gebaseerde opstart}

In plaats van een standaard VirtualBox-image wordt de virtuele schijf rechtstreeks gekoppeld via een VDI-bestand. Dit kan de opstart en provisioning vereenvoudigen doordat er minder extra conversie- of opbouwstappen nodig zijn. Het importeren van een volledige VM kan lang duren, terwijl het koppelen van een bestaande VDI dit proces mogelijk versnelt. Daardoor wordt de omgeving sneller beschikbaar voor gebruik.

Een belangrijk ontwerpdoel is het verkorten van de provisioningtijd. Door een VDI rechtstreeks te koppelen in plaats van met een volledig opgebouwde VirtualBox-image te werken, kan de opstartprocedure mogelijk sneller verlopen. Dit moet later bevestigd worden door een vergelijking tussen verschillende configuraties. De snelheid van dit proces is belangrijk omdat de gebruiker zo minder lang hoeft te wachten vooraleer de omgeving klaar is.




\subsubsection{VM's eerst opstarten}

Bij het gebruik van meerdere virtuele machines worden eerst alle VM's opgestart voordat de scripts worden uitgevoerd. Hierdoor kan de totale wachttijd verminderd worden, omdat de opstarttijd van de virtuele machine en de SSH-beschikbaarheid niet per machine apart moet worden afgewacht tijdens de provisioning.


\subsubsection{Netwerk- en opslagopzet}

Naast de VDI-keuze speelt ook de netwerkconfiguratie een rol in de totale opstarttijd. Door de VM meteen via een bridge adapter of vergelijkbare netwerkopzet te verbinden, kan het systeem sneller bereikbaar worden voor verdere provisioning. Dit vermindert de tijd tussen het opstarten van de VM en het moment waarop de provisioner effectief kan beginnen.



\subsubsection{Preconfiguratie van boxes}
Tijdens het opstarten van de virtuele machines wordt na het aanmaken en starten geen extra configuratie uitgevoerd. Bij Vagrant worden daarentegen nog verschillende eigenschappen van de nieuwe VM aangepast, zoals de hostname en de SSH-sleutel. Deze bijkomende stappen verlengen het opstartproces.




\subsection{Provisioningworkflow}

De provisioner doorloopt een vaste reeks stappen: de VM wordt opgestart, het IP-adres wordt opgehaald, er wordt verbinding gemaakt via SSH en vervolgens worden de provisioningstaken uitgevoerd. Deze reeks gebeurt in een vaste volgorde omdat elke stap afhankelijk is van de vorige. Zonder een correct opgestarte VM kan bijvoorbeeld geen netwerkverbinding worden gemaakt, en zonder SSH kan geen verdere configuratie gebeuren.

\begin{enumerate}
    \item \textbf{Aanmaken VM}
    
    In deze stap wordt de virtuele machine voorbereid en geregistreerd in de virtualisatieomgeving. Dit vormt de basis van het volledige provisioningproces. Daarna kunnen de netwerk- en opslaginstellingen worden aangepast zodat de machine correct kan starten.
    
    \item \textbf{Kopiëren van de VDI en aanpassing van de SID}
    
    Hier wordt de virtuele schijf gekopieerd of hergebruikt en worden de nodige instellingen aangepast zodat de VM correct kan opstarten. Deze stap zorgt ervoor dat de nieuwe omgeving een eigen identiteit krijgt en niet conflicteert met andere virtuele machines. Ook kunnen hier eventueel machine-specifieke instellingen worden ingesteld.
    
    \item \textbf{Bridgeadapter}
    
    Er wordt gebruikgemaakt van een bridgeadapter, zodat een tweede netwerkinterface aan de VM wordt toegevoegd. Deze keuze werd gemaakt omwille van de eenvoud. Het toevoegen van extra netwerkadapters zou de opzet complexer maken en extra elementen introduceren die voor deze benchmarks niet nodig zijn.
    
    \item \textbf{Shared folders}
    
    In deze fase wordt de \texttt{./script}-map gemount aan de virtuele machine. Hierdoor worden bestanden en scripts van de host direct toegankelijk binnen de VM. De shared folder wordt gebruikt zodat de provisioner, eenmaal verbonden via SSH, de benodigde scripts kan uitvoeren die zich in deze map bevinden.
    
    \item \textbf{NAT Port Forwarding}
    
    De netwerkconfiguratie wordt aangepast zodat de host toegang krijgt tot de virtuele machine via port forwarding. Dit is nodig om later via SSH verbinding te kunnen maken. Zonder deze stap zou de provisioner geen betrouwbare communicatie met de VM kunnen opzetten.
    
    \item \textbf{Opstarten van de VM}
    
    Na de configuratie wordt de virtuele machine effectief opgestart. Op dit moment wordt gecontroleerd of de VM correct reageert en of de benodigde netwerkparameters beschikbaar zijn. Deze stap is essentieel voordat de provisioning verder kan gaan.
    
    \item \textbf{SSH en shared folder}
    
    Zodra de VM beschikbaar is, wordt via SSH verbinding gemaakt en kunnen gedeelde mappen of andere provisioningstaken worden uitgevoerd. Via SSH kunnen configuratiescripts of installatieopdrachten op afstand worden uitgevoerd. Shared folders kunnen daarnaast gebruikt worden om bestanden tussen host en VM uit te wisselen.
\end{enumerate}


\subsubsection{Linux provisioning}


\begin{listing}[H]
    \begin{minted}{rust}
        channel.exec(&format!("echo {} | sudo -S bash /media/sf_shared/{}", password, scriptname))
        .expect("Command execution failed");
    \end{minted}
    \caption[Provisioning in Linux]{Uitvoering van het provisioningscript in Linux}
    \label{lst:linux-prov-line}
\end{listing}


\subsubsection{Windows provisioning}


\begin{listing}[H]
    \begin{minted}{rust}
        channel.exec(&format!(".\\psexec.exe -s -accepteula powershell.exe -ExecutionPolicy Bypass -File \"Z:\\{}\"", scriptname))
        .expect("Command execution failed ");
    \end{minted}
    \caption[Provisioning in Windows]{Uitvoering van het provisioningscript in Windows via psexec}
    \label{lst:windows-prov-line}
\end{listing}


\chapter{Uitbreidingen}
\label{ch:Uitbreidingen}

Naast het opstarten en provisionen voorziet de tool ook extra functies die het gebruik in labo-omgevingen verbeteren, zoals snapshots, foutafhandeling en logging. Deze functies maken het systeem robuuster en praktischer voor dagelijks gebruik. Ze zorgen er ook voor dat de gebruiker sneller kan herstellen wanneer er iets misloopt.

\subsection{Snapshotting}

Voor het provisionen kan een snapshot worden gemaakt, zodat de gebruiker bij fouten altijd kan terugkeren naar een stabiele toestand. Dit verhoogt de betrouwbaarheid, maar kost wel extra opslagruimte. In een onderwijscontext kan dit nuttig zijn om snel opnieuw te beginnen zonder alles opnieuw te moeten installeren.

\subsection{Priority Queue}

De VM’s kunnen in een vaste volgorde worden opgestart via een prioriteitsmechanisme. Die volgorde wordt als extra parameter in het configuratiebestand opgenomen. Dit kan handig zijn wanneer bepaalde machines afhankelijk zijn van andere machines die eerst beschikbaar moeten zijn.

\subsection{Kill}
Deze functie zorgt ervoor dat VirtualBox geforceerd wordt afgesloten. Ze werd toegevoegd omdat VirtualBox in sommige gevallen niet correct reageert bij het beheren van virtuele machines, bijvoorbeeld wanneer het programma vastloopt of foutcodes blijft geven.


\subsection{Ondersteuning}

\subsubsection{Ondersteunde besturingssystemen}

De code past zich automatisch aan op basis van het besturingssysteem waarop de provisioner wordt uitgevoerd. Hierdoor kan de juiste gebruikersmap of tijdelijke opslaglocatie worden gebruikt, afhankelijk van of de tool op Windows of Linux draait.

\begin{listing}[H]
    \begin{minted}{rust}
       pub fn get_home_dir() -> std::path::PathBuf {
           std::env::var("HOME")
           .map(std::path::PathBuf::from)
           .or_else(|_| {
               std::env::var("USERPROFILE")
               .map(std::path::PathBuf::from)
               .map_err(|_| std::io::Error::new(std::io::ErrorKind::NotFound, "Home directory not found"))
           })
           .expect("Failed to determine home directory")
       }
    \end{minted}
    \caption[Zoeken naar de gebruikersmap per besturingssysteem]{Voorbeeld van hoe de doelmap dynamisch wordt opgebouwd op basis van het besturingssysteem.}
    \label{lst:homedir-crossplatform}
\end{listing}

\begin{listing}[H]
    \begin{minted}{rust}
        if cfg!(target_os = "windows") {
            //code fragement
        }
    \end{minted}
    \caption[Zoeken naar het os type]{Zoeken naar het os type}
    \label{lst:ostype-crossplatform}
\end{listing}

\subsubsection{Ondersteunde besturingssystemen voor provisioning}

Er werd gekozen voor ondersteuning van Linux en Windows, omdat dit de twee besturingssysteemtypes zijn die binnen HOGENT het vaakst gebruikt worden. Voor deze bachelorproef werden vier boxes voorzien.

\paragraph{AlmaLinux}
AlmaLinux is de gekozen Linuxdistributie binnen HOGENT voor het vak Linux. Daarom wordt deze distributie ook ondersteund binnen de provisioner.

\paragraph{Ubuntu}
Ubuntu is een van de meest gebruikte Linuxdistributies en werd daarom  opgenomen als ondersteunde omgeving.
//https://www.openlogic.com/blog/top-enterprise-linux-distributions


\paragraph{Windows 10}
Windows 10 werd gekozen als clientbesturingssysteem voor de Windowsvakken.

\paragraph{Windows Server 2025}
Windows Server 2025 werd gekozen als serverbesturingssysteem voor de Windowsvakken.


één voor AlmaLinux, één voor Ubuntu en twee voor Windows-omgevingen. De Windows boxes zijn voor Windows 10 en voor Windows Server 2025

\subsection{Dependencies}

Voor de implementatie worden onder andere de volgende Rust-crates gebruikt:
\begin{itemize}
    \item \texttt{serde} voor het serialiseren en deserialiseren van gegevens. Deze crate werd gebruikt in combinatie met \texttt{toml} om configuratiebestanden vlot om te zetten naar Rust-structuren en omgekeerd.
    
    \item \texttt{ssh2} voor SSH-verbindingen. Deze crate werd gekozen voor cross-platform compatibiliteit, omdat de ingebouwde SSH-oplossingen van Linux en Windows te verschillend zijn en niet op dezelfde manier met authenticatie omgaan. De crate is een Rust-binding voor \texttt{libssh2}, een SSH-clientbibliotheek, en is dual-licensed onder Apache 2.0 en MIT. \autocite{ssh2Crate}
    
    \item \texttt{toml} voor het inlezen van configuratiebestanden. Deze crate werd gekozen omwille van de eenvoud en duidelijkheid van het importeren van data uit een TOML-bestand. TOML is een leesbaar configuratieformaat dat vlot kan worden omgezet naar Rust-structuren. De crate is een serde-compatibele TOML-decoder en -encoder voor Rust en is dual-licensed onder Apache 2.0 en MIT. \autocite{tomlCrate}
    
    \item \texttt{clap} voor command-line argumenten. Deze crate werd gekozen omdat ze een eenvoudige en gebruiksvriendelijke manier biedt om CLI-argumenten te verwerken. De crate is dual-licensed onder Apache 2.0 en MIT. \autocite{clapCrate}
    
    \item \texttt{indexmap} voor het bewaren van invoervolgorde bij configuratiegegevens. Deze crate werd gebruikt omdat de volgorde van de VM's belangrijk is bij het opstarten en provisionen, iets wat met een standaard \texttt{HashMap} niet gegarandeerd wordt.
    
    \item \texttt{crossterm} voor terminalinteractie en console-uitvoer. Deze crate werd gebruikt om de command-line ervaring overzichtelijker en gebruiksvriendelijker te maken.
    
    \item \texttt{PSTools} voor Windows-specifieke uitvoering van scripts als SYSTEM-gebruiker.
\end{itemize}



\chapter{Benchmark}
\label{ch:Benchmark}

\subsection{Snelheidsvergelijking met Vagrant}

De uiteindelijke implementatie wordt vergeleken met Vagrant om na te gaan of de nieuwe provisioner sneller of trager is. Hiervoor wordt de provisioningtijd gemeten onder gelijkaardige omstandigheden. De vergelijking gebeurt op basis van een beperkt aantal testgevallen, namelijk drie Linux-VM’s en twee Windows-VM’s. Op die manier kan worden nagegaan of de gekozen aanpak een echte verbetering vormt ten opzichte van de bestaande oplossing.

Deze testen werden uitgevoerd op een computer met 32 GB DDR4-geheugen aan 3200 MHz, een AMD Ryzen 7 5800H-processor en een NVIDIA GeForce RTX 3060-laptopgpu met een vermogen van 90 W op Windows 11 25H2.

\begin{itemize}
    \item \texttt{Eén VM: } Tijd vanaf het opstarten van de VM tot de start van het provisioningproces. Deze meting wordt stopgezet op het moment dat de provisioning begint, omdat dit een beter beeld geeft van de snelheid van de provisioner zelf. Het uitvoeren van het provisioningscript wordt namelijk mede beïnvloed door externe factoren zoals internetverbindingen, wat de meting minder stabiel maakt.
    
    \item \texttt{Meerdere VM's: } Totale uitvoertijd van alle VM's, na elkaar uitgevoerd. Deze meting is realistischer omdat ook de uitvoering van de scripts wordt meegerekend, waardoor een meer volledige vergelijking met meerdere VM's mogelijk is.
\end{itemize}


\subsubsection{Windows}

Voor Windows worden een Windows Server 2025-omgeving en een Windows 10-clientomgeving getest. De Windows-gevallen zijn gebaseerd op de SEP-opdracht van vorig jaar en bevatten een domeincontroller en een client die via DHCP wordt toegevoegd aan het domein. Eerst wordt de tijd gemeten tot de provisioning van één VM volledig afgerond is. Daarna wordt ook de totale tijd gemeten wanneer meerdere VM’s worden opgestart en alle scripts volledig zijn uitgevoerd.

\begin{table}[H]
    \centering
    \begin{tabular}{lcccc}
        \toprule
        \textbf{Test} & \textbf{1} & \textbf{2} & \textbf{3} & \textbf{Gemiddelde} \\
        \midrule
        Eén VM: Vagrant & x & x & x & x \\
        Eén VM: RustProv & x & x & x & x \\
        Meerdere VM's: Vagrant & x & x & x & x \\
        Meerdere VM's: RustProv & x & x & x & x \\
        \bottomrule
    \end{tabular}
    \caption{Vergelijking van de Windows-tijden tussen Vagrant en RustProv in seconden.}
    \label{tab:windows-vergelijking}
\end{table}

\subsubsection{Linux}

Voor Linux worden drie virtuele machines op een gelijkaardige manier geanalyseerd. Ook hier wordt gekeken naar de tijd tot de provisioning van één VM volledig afgerond is. Vervolgens wordt de totale uitvoeringstijd gemeten voor het scenario met meerdere VM’s, waarbij alle scripts uitgevoerd en afgerond zijn-. Op die manier kan een eerlijke vergelijking gemaakt worden tussen Vagrant en RustProv.

\begin{table}[H]
    \centering
    \begin{tabular}{lcccc}
        \toprule
        \textbf{Test} & \textbf{1} & \textbf{2} & \textbf{3} & \textbf{Gemiddelde} \\
        \midrule
        Eén VM: Vagrant & x & x & x & x \\
        Eén VM: RustProv & x & x & x & x \\
        Meerdere VM's: Vagrant & x & x & x & x \\
        Meerdere VM's: RustProv & x & x & x & x \\
        \bottomrule
    \end{tabular}
    \caption{Vergelijking van de Linux-tijden tussen Vagrant en RustProv in seconden.}
    \label{tab:linux-vergelijking}
\end{table}